\chapter{Vorbereitung}

    \section{Kräfte zwischen Atomen}
        \subsection{Van-der-Waals Kräfte}

\begin{figure}[htb]
    \centering
    \includegraphics[width=0.6\textwidth]{Abb/vdw.pdf}
    \caption{Van-der-Waals Kraft zwischen zwei Teilchen}
    \label{vdw}
\end{figure}

Die Ladungsverteilung in Atomen ist nicht konstant, sondern unterliegt ständiger 
Fluktuation. Der Schwerpunkt der negativen Ladungen kann hierbei vom dem der
positiven Ladungen abweichen. Ist dies der Fall, so entsteht ein Dipol. 
Befindet sich nun ein zweites Teilchen in der Nähe dieses Atoms, so wird auch
in diesem ein Dipol induziert. Zeigt die positive Seite des ersten Atoms zu Atom 2,
so werden die Elektronen des zweiten Atoms angezogen. Ist es die negative Seite, so
werden die Elektronen abgestoßen. \par
Als Folge dessen synchronisieren sich die Ladungsänderungen der beiden Atome. Eine
schwache positive Anziehung ist die Folge. Diese ist proportional zu $\displaystyle
\frac{-1}{r^6}$.

        \subsection{Pauli-Abstoßung}

Nähern sich die Atome weiter an, so kommt es zu einem Überlappen der 
Elektronenorbitale. Das Pauli-Verbot verhindert hierbei, dass zwei Elektronen den
gleichen Zustand besetzen. Einige Elektronen werden folglich in einen energetisch
höheren Zustand gezwungen. \par
So führt eine Orbitalüberlagerung zu einer repulsiven Wechselwirkung. Die Kraft
ist proportional zu $\displaystyle \frac{1}{r^{12}}$.

        \subsection{Lennard-Jones Potential}

Bei sehr kleinen Abständen dominiert die Pauli-Abstoßung, bei größeren die 
van-der-Waals Kräfte. Die Summe aus beiden Potentialen wird Lennard-Jones 
Potential gennant. 
\[
   \phi (r) \propto \frac{A}{r^6} - \frac{B}{r^{12}}    
\]

\begin{figure}[hb]
    \centering
    \includegraphics[width=0.6\textwidth]{Abb/ljp.jpg}
    \caption{Das Lennard-Jones Potential als Summe der vdW-Wechselwirkung und
             der Pauli-Abstoßung}
    \label{ljp}
\end{figure}
